\documentclass[10pt,a4paper]{amsart}
\usepackage[margin=19mm]{geometry}
\usepackage{amsmath,amssymb,mathtools}
\usepackage[scr=rsfs,cal=euler]{mathalfa}
\usepackage{xfrac}
\usepackage[unicode,hidelinks,bookmarks=false]{hyperref}
\newcommand{\ie}{i.e.\ }
\newcommand{\cf}{cf.\ }
\newcommand{\resp}{respectively\ }
\newcommand{\Subsection}{\subsection}
\providecommand{\nobreakdash}{\nobreak\hbox{-}}
\setlength{\emergencystretch}{2em}
\multlinegap0pt
\usepackage{mdframed}
\usepackage{mdframed}
\usepackage{mdframed}
\usepackage{mdframed}
\usepackage{algorithm}\usepackage{algpseudocode}
\usepackage{bm}
\usepackage{xcolor}
\usepackage{tikz}
\usepackage{amssymb}
\usepackage[shortlabels]{enumitem}
\usepackage{tcolorbox}
\usepackage{mathtools}
\usepackage{mdframed}
\usepackage[normalem]{ulem}
\newtheorem{proposition}{Proposition}
\usepackage{cleveref}
\crefname{proposition}{Proposition}{Propositions}
\providecommand{\paragraph}{}
\title{Select-then-differentiate: Solving Bilevel Optimization with Manifold Lower-level Solution Sets}
\author{Saeed Masiha, Zebang Shen, Negar Kiyavash, Niao He}
\date{Source: arXiv:2605.09209v1 (CC BY 4.0)}
\begin{document}
\maketitle

\noindent\emph{Source and licence.} Select-then-differentiate: Solving Bilevel Optimization with Manifold Lower-level Solution Sets. Version arXiv:2605.09209v1 (CC BY 4.0), distributed by arXiv under the Creative Commons Attribution 4.0 International licence, http://creativecommons.org/licenses/by/4.0/, as recorded in the arXiv licence metadata for this exact version and shown on the abstract page https://arxiv.org/abs/2605.09209v1. Full licence notice: \texttt{licenses/arxiv-2605.09209-CC-BY-4.0.txt}.\par\smallskip

\noindent\textit{Editorial scope.} Counted unit: the article's proposition prop:outer-inexact (source0.tex lines 1373--1384). The article's complete original proof of it is delivered with it (source0.tex lines 3034--3105, one \texttt{proof} environment of 72 lines, which the article places away from the statement). No mathematics is written, repaired or re-derived: the statement and the proof are the article's own lines, in the article's own notation, and no retained line is substituted or re-flowed.  Retained uncounted as the source-local context the proof needs: lines 1287--1339.  Nothing was omitted from any retained span, so the package carries no omission record.  No retained line contains a citation, so the wrapper carries no bibliography and no bibliography entry.  No retained line contains a figure environment, an image-inclusion command or drawing code, so no image is delivered and the package declares no asset dependency.  Editorial formatting only, in addition: an AMS \texttt{amsart} preamble that restates the article's own declarations and macros used by the retained lines, namely the environments \texttt{proposition} on separate counters, together with the standard packages those lines need.  The retained environments print in this document as Proposition 1 in order of appearance, whereas the article prints the same items with the numbering of its own full document.  The complete native source of the article as distributed in the arXiv e-print for this exact version is bundled unchanged as \texttt{sources/200206/source0.tex}.
\begin{algorithm}[t]
\caption{\textsc{HG-MS}: Hyper-gradient with minima selection}
\begin{algorithmic}[1]
\Require initial $\theta_0\in\Theta$; stepsizes $\{\alpha_t\}$; temperature $\lambda$; number of chains $N$; ULA steps $K$ and stepsize $h$; chain starts $\{X^{(0)}_{t,i}\}$; CG tolerances $\{\eta_t\}$; ridge $\gamma\ge 0$; clip radius $R_v>0$
\For{$t=0,1,2,\ldots$}
    \State Sample in parallel $X_{t,i}\gets \textsc{ULA}(\theta_t,g,X^{(0)}_{t,i},\lambda,K,h)$
    % \For{$i=1$ to $N$ \textbf{in parallel}}
    %     \State 
    % \EndFor
    \State $i_t^\star\gets\arg\min_{1\le i\le N} f(\theta_t,X_{t,i})$; \quad
    $\hat{x}_N^\lambda(\theta_t)\gets X_{t,i_t^\star}$
    \State Solve
    \(
        \big(\nabla^2_{xx}g(\theta_t,\hat{x}_N^\lambda(\theta_t))+\gamma I\big)\tilde v_t
        =
        \nabla_x f(\theta_t,\hat{x}_N^\lambda(\theta_t))
    \)
    by CG up to residual $\eta_t$
    \State \textbf{(Optional safeguard)} $v_t\gets \mathrm{Proj}_{\mathbb B_d(0;R_v)}(\tilde v_t)$
    \State $\widehat h_t\gets
    \nabla_\theta f(\theta_t,\hat{x}_N^\lambda(\theta_t))
    -
    \nabla^2_{\theta x}g(\theta_t,\hat{x}_N^\lambda(\theta_t))v_t$; $\quad \theta_{t+1}\gets
    \mathrm{Proj}_{\Theta}\big(\theta_t-\alpha_t\widehat h_t\big)$
\EndFor
\end{algorithmic}
% \begin{algorithmic}[1]
% \INPUT initial $\theta_0\in\Theta$; stepsizes $\{\alpha_t\}$; temperature $\lambda$; number of chains $N$; ULA steps $K$ and stepsize $h$; CG tolerances $\{\eta_t\}$; ridge $\gamma\ge 0$; clip radius $R_v>0$.
% \STATE Initialize chain states $X^{(0)}_{0,1},\ldots,X^{(0)}_{0,N}$.
% \FOR{$t=0,1,2,\ldots$}
%     \FOR{$i=1$ to $N$ \textbf{in parallel}}
%         \STATE $X_{t,i}\leftarrow \textsc{ULA}(\theta_t,g,X^{(0)}_{t,i},\lambda,K,h)$.
%     \ENDFOR
%     \STATE $i_t^\star\leftarrow\arg\min_{1\le i\le N} f(\theta_t,X_{t,i})$; \quad
%     $\hat{x}_N^\lambda(\theta_t)\leftarrow X_{t,i_t^\star}$.
%     \STATE Solve
%     \(
%         \big(\nabla^2_{xx}g(\theta_t,\hat{x}_N^\lambda(\theta_t))+\gamma I\big)\tilde v_t
%         =
%         \nabla_x f(\theta_t,\hat{x}_N^\lambda(\theta_t))
%     \)
%     by CG up to residual $\eta_t$.
%     \STATE \textbf{(Safeguard)} $v_t\leftarrow \mathrm{Proj}_{\mathbb B_d(0;R_v)}(\tilde v_t)$; otherwise set $v_t=\tilde v_t$.
%     \STATE $\widehat h_t\leftarrow
%     \nabla_\theta f(\theta_t,\hat{x}_N^\lambda(\theta_t))
%     -
%     \nabla^2_{\theta x}g(\theta_t,\hat{x}_N^\lambda(\theta_t))v_t$.
%     \STATE $\theta_{t+1}\leftarrow
%     \mathrm{Proj}_{\Theta}\big(\theta_t-\alpha_t\widehat h_t\big)$.
% \ENDFOR
% \end{algorithmic}
\label{alg:hg-minsel-lmc}
\end{algorithm}

\begin{proposition}[Convergence to near-stationarity with inexact hyper-gradients]\label{prop:outer-inexact}
Suppose that $F$ is $L_F$-smooth on $\Theta$.
Let $F_\star:=\inf_{\theta\in\Theta}F(\theta)>-\infty$.
Initialize $\theta_0\in\Theta$ and run the projected update
$\theta_{t+1}=\mathrm{Proj}_{\Theta}(\theta_t-\alpha \widehat h_t)$ with a constant
stepsize $\alpha\le 1/L_F$ (last line of \Cref{alg:hg-minsel-lmc}), then
\begin{align}
\forall\ T \geq 1, \quad \frac{1}{T}\sum_{t=0}^{T-1}\mathbb E\big[\|\mathcal G_{\Theta}(\theta_t,\nabla F(\theta_t);\alpha)\|^2\big]\le
\frac{8\,(F(\theta_0)-F_\star)}{\alpha T}+ \frac{10}{T}\sum_{t=0}^{T-1}\mathbb E\big[\|e_t\|^2\big].
\label{eq:outer-inexact}
\end{align}
\end{proposition}

\begin{proof}[Proof of \Cref{prop:outer-inexact}]
Fix $T\ge 1$ and write $g_t:=\widehat h_t=\nabla F(\theta_t)+e_t$.
Define the (used) gradient mapping
\[
\widetilde{\mathcal G}_t
:=\alpha^{-1}\big(\theta_t-\mathrm{Proj}_{\Theta}(\theta_t-\alpha g_t)\big)
=\alpha^{-1}(\theta_t-\theta_{t+1}).
\]

\paragraph{Step 1: a descent inequality in terms of $\widetilde{\mathcal G}_t$.}
Let $\Delta_t:=\theta_{t+1}-\theta_t=-\alpha\,\widetilde{\mathcal G}_t$.
By $L_F$-smoothness of $F$ on $\Theta$ and convexity of $\Theta$, we have
\[
F(\theta_{t+1})
\le
F(\theta_t)+\langle \nabla F(\theta_t),\Delta_t\rangle+\frac{L_F}{2}\|\Delta_t\|^2.
\]
Since $\theta_{t+1}=\mathrm{Proj}_{\Theta}(\theta_t-\alpha g_t)$ and $\theta_t\in\Theta$, the projection
optimality condition yields
\[
\langle g_t,\Delta_t\rangle\ \le\ -\frac{1}{\alpha}\|\Delta_t\|^2.
\]
	Using $\nabla F(\theta_t)=g_t-e_t$, Cauchy--Schwarz, and the weighted Young inequality
	\[
	uv\le \frac{1}{4\alpha}u^2+\alpha v^2,\qquad u,v\ge0,
	\]
	gives
\begin{align*}
F(\theta_{t+1})
&\le
F(\theta_t)+\langle g_t,\Delta_t\rangle-\langle e_t,\Delta_t\rangle+\frac{L_F}{2}\|\Delta_t\|^2\\
&\le
F(\theta_t)-\frac{1}{\alpha}\|\Delta_t\|^2+\|e_t\|\,\|\Delta_t\|+\frac{L_F}{2}\|\Delta_t\|^2\\
&\le
F(\theta_t)-\frac{1}{\alpha}\|\Delta_t\|^2+\frac{1}{4\alpha}\|\Delta_t\|^2+\alpha\|e_t\|^2+\frac{L_F}{2}\|\Delta_t\|^2\\
&=
F(\theta_t)-\Big(\frac{3}{4\alpha}-\frac{L_F}{2}\Big)\|\Delta_t\|^2+\alpha\|e_t\|^2.
\end{align*}
Using $\alpha\le 1/L_F$ gives $\frac{3}{4\alpha}-\frac{L_F}{2}\ge \frac{1}{4\alpha}$, hence
\[
F(\theta_{t+1})
\le
F(\theta_t)-\frac{1}{4\alpha}\|\Delta_t\|^2+\alpha\|e_t\|^2
=
F(\theta_t)-\frac{\alpha}{4}\|\widetilde{\mathcal G}_t\|^2+\alpha\|e_t\|^2.
\]
Summing over $t=0,\dots,T-1$ and using $F(\theta_T)\ge F_\star$ yields
\begin{equation}\label{eq:sum-Gtilde}
\sum_{t=0}^{T-1}\|\widetilde{\mathcal G}_t\|^2
\le
\frac{4\,(F(\theta_0)-F_\star)}{\alpha}
+4\sum_{t=0}^{T-1}\|e_t\|^2.
\end{equation}

\paragraph{Step 2: relate $\widetilde{\mathcal G}_t$ to the true gradient mapping.}
By nonexpansiveness of projection,
\[
\|\mathcal G_{\Theta}(\theta_t,\nabla F(\theta_t);\alpha)-\widetilde{\mathcal G}_t\|
=
\frac{1}{\alpha}\big\|\mathrm{Proj}_{\Theta}(\theta_t-\alpha g_t)-\mathrm{Proj}_{\Theta}(\theta_t-\alpha\nabla F(\theta_t))\big\|
\le \|e_t\|.
\]
Therefore $\|\mathcal G_{\Theta}(\theta_t,\nabla F(\theta_t);\alpha)\|^2\le 2\|\widetilde{\mathcal G}_t\|^2+2\|e_t\|^2$,
and combining with \eqref{eq:sum-Gtilde} gives
\[
\sum_{t=0}^{T-1}\|\mathcal G_{\Theta}(\theta_t,\nabla F(\theta_t);\alpha)\|^2
\le
\frac{8\,(F(\theta_0)-F_\star)}{\alpha}
+10\sum_{t=0}^{T-1}\|e_t\|^2.
\]
Dividing by $T$ and taking expectations yields \eqref{eq:outer-inexact}.
\end{proof}

\end{document}
